\chapter{সাহিত্য পর্যালোচনা}\label{ch:literature}

এই অধ্যায় প্রকল্পটিকে বিদ্যমান জ্ঞানভাণ্ডারের প্রেক্ষাপটে স্থাপন করে। অ্যানালিটিক্সের ব্যবসায়িক যৌক্তিকতা থেকে শুরু করে বিশ্লেষণী
ব্যবস্থা নির্মাণের প্রক্রিয়া মডেল ও কৌশল, এবং এই কাজে ব্যবহৃত নির্দিষ্ট পদ্ধতিগুলো আলোচনা করে শেষে প্রকল্পটি যে গবেষণা-ঘাটতি
পূরণ করে তা চিহ্নিত করা হয়েছে।

\section{বিজনেস ইন্টেলিজেন্স ও অ্যানালিটিক্স}
\citet{chen2012} বিজনেস ইন্টেলিজেন্স ও অ্যানালিটিক্সের (BI\&A) বিবর্তনকে তিন পর্যায়ে বর্ণনা করেছেন। প্রথম পর্যায়, BI\&A~1.0,
ডেটা ওয়্যারহাউস ও কাঠামোবদ্ধ প্রতিবেদনকেন্দ্রিক। দ্বিতীয় পর্যায়ে যুক্ত হয় ওয়েব ও অসংগঠিত ডেটা, এবং তৃতীয় পর্যায়ে মোবাইল ও সেন্সর
ডেটা। তাঁদের মতে মূল মূল্য নিহিত ডেটাকে \emph{সিদ্ধান্তে} রূপান্তরে, কেবল প্রতিবেদনে নয়। \citet{davenport2007} প্রাতিষ্ঠানিক কেস
স্টাডি থেকে একই সিদ্ধান্তে পৌঁছেছেন: যেসব প্রতিষ্ঠান ``অ্যানালিটিক্স দিয়ে প্রতিযোগিতা করে'', তারা মাঝে মাঝে প্রকল্পে নয়, বরং মূল
কার্যপ্রক্রিয়ায় নিয়মিতভাবে পরিসংখ্যানিক ও পূর্বাভাসমূলক মডেল ব্যবহার করে। \citet{sharda2018} অ্যানালিটিক্সকে তিন স্তরে
সাজিয়েছেন: \emph{বর্ণনামূলক} (কী ঘটেছে), \emph{পূর্বাভাসমূলক} (কী ঘটবে) এবং \emph{নির্দেশমূলক} (কী করা উচিত)। অধিকাংশ প্রতিষ্ঠান
এখনো বর্ণনামূলক স্তরে আছে। প্রকল্প প্রস্তাবনার তিনটি ব্যবস্থাপনা প্রশ্ন সরাসরি এই স্তরবিন্যাসের সঙ্গে মেলে; এখানে নির্মিত কাঠামো
বর্ণনামূলক ও পূর্বাভাসমূলক স্তর আচ্ছাদন করে এবং নিয়ম-ভিত্তিক নির্দেশনা দেয়।

\citet{provost2013} \emph{ডেটা-বিশ্লেষণী চিন্তার} উপর জোর দিয়েছেন: অ্যানালিটিক্স তখনই মূল্যবান যখন ব্যবসায়িক সমস্যাটি এতটা
সুনির্দিষ্টভাবে সংজ্ঞায়িত হয় যে ডেটা দিয়ে উত্তর দেওয়া যায়, এবং মডেলের বিচার হওয়া উচিত সিদ্ধান্তে তার প্রভাব দিয়ে, কেবল প্রযুক্তিগত
নির্ভুলতা দিয়ে নয়। \citet{shmueli2010} যোগ করেছেন যে ব্যাখ্যামূলক ও পূর্বাভাসমূলক মডেলিং দুটি ভিন্ন কাজ, যাদের লক্ষ্য ও যাচাই
পদ্ধতি আলাদা। যে মডেল অতীতের পরিবর্তন ভালো ব্যাখ্যা করে, সেটি দুর্বল পূর্বাভাস দিতে পারে, এবং উল্টোটাও সত্য। এই পার্থক্য
বর্তমান গবেষণাকে প্রভাবিত করেছে: ব্যাখ্যামূলক বিশ্লেষণ (EDA, হাইপোথিসিস টেস্ট, elasticity রিগ্রেশন) পূর্বাভাসমূলক বিশ্লেষণ
(out-of-sample পূর্বাভাস ও ক্লাসিফিকেশন) থেকে আলাদা রাখা হয়েছে এবং প্রতিটিকে তার নিজস্ব মানদণ্ডে মূল্যায়ন করা হয়েছে।

\section{অ্যানালিটিক্স প্রকল্পের প্রক্রিয়া মডেল}
অ্যানালিটিক্স প্রকল্পে সবচেয়ে বেশি ব্যবহৃত পদ্ধতি এখনো Cross-Industry Standard Process for Data Mining (CRISP-DM)
\citep{shearer2000}। এর ছয়টি ধাপ: ব্যবসা বোঝা, ডেটা বোঝা, ডেটা প্রস্তুতি, মডেলিং, মূল্যায়ন ও বাস্তবায়ন। ধাপগুলো পুনরাবৃত্তিমূলক,
এবং মডেলটি স্বীকার করে যে সাধারণত ডেটা প্রস্তুতিতেই সবচেয়ে বেশি শ্রম লাগে। প্রকল্প প্রস্তাবনার দশ-ধাপের পরিকল্পনা CRISP-DM-এর
একটি পরিমার্জিত রূপ, যেখানে ফিচার ইঞ্জিনিয়ারিং, ড্যাশবোর্ড উন্নয়ন, সিস্টেম মূল্যায়ন ও ডকুমেন্টেশনের জন্য আলাদা ধাপ রয়েছে।
\citet{han2012} একই রকম জ্ঞান-আবিষ্কার প্রক্রিয়া বর্ণনা করেছেন (পরিষ্কারকরণ, একীকরণ, নির্বাচন, রূপান্তর, মাইনিং, মূল্যায়ন ও
উপস্থাপন), এবং \cref{ch:system}-এর কাঠামো উভয়টিই অনুসরণ করে।

\section{ডেটার মান ও প্রিপ্রসেসিং}
বাস্তব ডেটা খুব কমই বিশ্লেষণের জন্য প্রস্তুত থাকে। \citet{garcia2015} প্রিপ্রসেসিংকে ভাগ করেছেন ডেটা পরিষ্কারকরণ (অনুপস্থিত মান,
noise ও outlier), একীকরণ, রূপান্তর (normalisation ও discretisation) এবং সংকোচনে। তাঁরা জোর দেন যে কৌশল নির্বাচন নির্ভর
করা উচিত সমস্যাটি কোন প্রক্রিয়ায় সৃষ্টি হয়েছে তার উপর। অনুপস্থিত মানের ক্ষেত্রে সারি মুছে ফেলা নিরাপদ কেবল তখনই, যখন মানগুলো
সম্পূর্ণ এলোমেলোভাবে অনুপস্থিত এবং সংখ্যায় কম। অন্যথায় imputation শ্রেয়, আর কলামের মধ্যকার নির্ধারিত সম্পর্ক থেকে মান
নির্ণয় করাই সবচেয়ে নির্ভুল পথ। \citet{wickham2014} যোগ করেছেন \emph{tidy data}-র নীতি: প্রতিটি চলক একটি কলাম, প্রতিটি
পর্যবেক্ষণ একটি সারি, এবং প্রতিটি ধরনের পর্যবেক্ষণ-একক একটি সারণি। এখানে ব্যবহৃত স্টার-স্কিমা নকশার ভিত্তি এই নীতি।

ব্যবসায়িক ডেটায় outlier সামলাতে বিশেষ সতর্কতা দরকার। $z$-score-এর মতো পরিসংখ্যানিক নিয়ম মোটামুটি স্বাভাবিক (normal) বণ্টন ধরে
নেয় এবং যে outlier খোঁজা হচ্ছে সেগুলো দিয়েই বিকৃত হয়। Inter-quartile range (IQR) নিয়ম ও median absolute deviation (MAD)
বেশি robust \citep{han2012}। তবে প্রতিটি outlier ভুল নয়: ডিস্ট্রিবিউটরের একটি বাল্ক অর্ডার অস্বাভাবিক কিন্তু প্রকৃত, এবং তা মুছে
ফেললে চাহিদার অনুমান নিচের দিকে পক্ষপাতদুষ্ট হবে। তাই এখানে \emph{অসম্ভব} মান (ডোমেইন-নিয়ম লঙ্ঘন, যা সংশোধন করা হয়) এবং
\emph{অস্বাভাবিক} মান (পরিসংখ্যানিক outlier, যা চিহ্নিত করে রেখে দেওয়া হয়) আলাদা করা হয়েছে।

আরেকটি প্রায়ই উপেক্ষিত ঝুঁকি হলো \emph{data leakage}: পূর্বাভাসের মুহূর্তে যে তথ্য পাওয়া যেত না, তা প্রশিক্ষণ ডেটায় ঢুকে পড়া।
\citet{kaufman2012} দেখিয়েছেন যে প্রকাশিত গবেষণা ও প্রতিযোগিতায় leakage সাধারণ ঘটনা, এবং এর ফলে মডেল পরীক্ষায় চমৎকার দেখায়
কিন্তু বাস্তব প্রয়োগে ব্যর্থ হয়। এর সাধারণ উৎস হলো লক্ষ্য-চলক থেকেই তৈরি ফিচার, বর্তমান সময়কাল অন্তর্ভুক্ত করা rolling পরিসংখ্যান,
এবং প্রশিক্ষণ অংশের বদলে পুরো ডেটাসেটে নির্ণীত প্রিপ্রসেসিং প্যারামিটার (যেমন scaling ধ্রুবক বা target encoding)। এই প্রকল্পে
নকশাগতভাবেই এই তিনটি এড়ানো হয়েছে।

\section{ডেটা ওয়্যারহাউজিং}
\citet{kimball2013} বিশ্লেষণী ডেটাবেসের মানক নকশা হিসেবে ডাইমেনশনাল মডেলিং প্রতিষ্ঠা করেছেন। একটি কেন্দ্রীয় \emph{fact table}
পরিমাপযোগ্য ঘটনা (এখানে বিক্রয় লেনদেন) সংরক্ষণ করে, এবং চারপাশের \emph{dimension table}-গুলো প্রতিটি ঘটনার প্রেক্ষাপট (তারিখ,
পণ্য, গ্রাহক বিভাগ) বর্ণনা করে। ফলস্বরূপ স্টার স্কিমা বোঝা সহজ, aggregation কোয়েরিতে দক্ষ, এবং Power BI-এর মতো BI টুলে সরাসরি
ব্যবহারযোগ্য। এই প্রকল্পে বহনযোগ্যতার জন্য SQLite-এ একটি স্টার স্কিমা তৈরি করা হয়েছে; একই SQL PostgreSQL বা MySQL-এও চলে।

\section{চাহিদা পূর্বাভাস}
\subsection{পরিসংখ্যানিক পদ্ধতি}
Exponential smoothing-এর সূচনা Holt ও \citet{winters1960}-এর কাজে, যিনি একে প্রবণতা ও মৌসুমি প্রভাবসহ সিরিজে সম্প্রসারিত
করেন। Holt--Winters পদ্ধতি ও এর আধুনিক state-space রূপ এখনো শক্তিশালী বেঞ্চমার্ক, যেমন seasonal ARIMA (SARIMA) পরিবার
\citep{hyndman2021}। উভয়ই একটি একক time series মডেল করে এবং ভালোভাবে ক্যালিব্রেটেড পূর্বাভাস-ব্যবধান দেয়, কিন্তু অতিরিক্ত
ব্যাখ্যাকারী চলক সহজে ব্যবহার করতে পারে না বা অনেক সম্পর্কিত সিরিজ থেকে যৌথ প্যাটার্ন শিখতে পারে না।

\subsection{মেশিন লার্নিং পদ্ধতি}
মেশিন লার্নিং রিগ্রেসর পূর্বাভাসকে supervised learning হিসেবে দেখে: পূর্ববর্তী (lagged) মান, rolling পরিসংখ্যান ও ক্যালেন্ডার
চলক হয়ে ওঠে ফিচার, এবং একটি সমন্বিত মডেল একসঙ্গে অনেক সিরিজে প্রশিক্ষিত হতে পারে। Tree ensemble সবচেয়ে প্রচলিত পছন্দ।
Random forest bootstrap নমুনায় প্রশিক্ষিত অনেক কম-সম্পর্কিত decision tree-র গড় নেয় \citep{breiman2001}। Gradient boosting
বর্তমান ensemble-এর residual-এর উপর ধারাবাহিকভাবে tree বসায় \citep{friedman2001}। XGBoost হলো gradient boosting-এর একটি
regularised ও অত্যন্ত অপ্টিমাইজড বাস্তবায়ন \citep{chen2016}। Support vector machine ও support vector regression একটি
kernel-সৃষ্ট ফিচার-স্পেসে সর্বোচ্চ-মার্জিন ফাংশন খোঁজে \citep{cortes1995}। $k$-nearest neighbours সবচেয়ে সদৃশ ঐতিহাসিক
উদাহরণ থেকে পূর্বাভাস দেয়। \citet{hastie2009} ও \citet{james2021} এসব পদ্ধতির তাত্ত্বিক ভিত্তি দেন, বিশেষত bias--variance
ভারসাম্য, যা regularisation ও hyper-parameter tuning-এর প্রেরণা।

\subsection{পরিসংখ্যানিক বনাম মেশিন লার্নিং পূর্বাভাস}
\citet{makridakis2018} 1,045টি মাসিক সিরিজে আটটি পরিসংখ্যানিক ও দশটি মেশিন লার্নিং পদ্ধতির তুলনা করে দেখান যে গড়ে সরল
পরিসংখ্যানিক পদ্ধতিগুলো বেশি নির্ভুল ছিল। তাঁদের উপসংহার: মেশিন লার্নিং পূর্বাভাসকে সবসময় সরল পদ্ধতির সঙ্গে তুলনা করতে হবে।
এর দুটি ফল এই প্রকল্পে প্রতিফলিত। প্রথমত, naive ও seasonal-naive পূর্বাভাস এবং Holt--Winters ও SARIMA বেঞ্চমার্ক হিসেবে
অন্তর্ভুক্ত। দ্বিতীয়ত, মূল্যায়ন নকশা পূর্বাভাসের সর্বোত্তম চর্চা অনুসরণ করে। Time-series ডেটা যাচাইয়ের জন্য এলোমেলোভাবে মিশানো
উচিত নয়, কারণ তাতে মডেল ``ভবিষ্যৎ দেখে ফেলে''; rolling-origin বা expanding-window পদ্ধতি শ্রেয় \citep{bergmeir2012,
hyndman2021}। বিভিন্ন স্তরের (পণ্য বনাম কোম্পানি) পূর্বাভাস reconciliation-এর মাধ্যমে সামঞ্জস্যপূর্ণ করা যায়
\citep{athanasopoulos2009}।

\subsection{নির্ভুলতার পরিমাপ}
\citet{hyndman2006} পূর্বাভাস-নির্ভুলতার পরিমাপগুলো পর্যালোচনা করেছেন। স্কেল-নির্ভর পরিমাপ (MAE, RMSE) ডেটার এককেই বোঝা
যায়। শতকরা ত্রুটি (MAPE) স্কেল-মুক্ত, কিন্তু প্রকৃত মান শূন্যের কাছাকাছি হলে বিস্ফোরিত হয়। বেঞ্চমার্ক-আপেক্ষিক পরিমাপ দেখায় একটি
পদ্ধতি আসলে মূল্য যোগ করছে কি না। তাই এই গবেষণায় MAE, MSE, RMSE, $R^2$, MAPE ও weighted absolute percentage error
(WAPE) এবং naive বেঞ্চমার্কের তুলনায় উন্নতি উল্লেখ করা হয়েছে।

\section{ক্লাসিফিকেশন ও মূল্যায়ন}
ক্লাসিফিকেশন একটি শ্রেণিগত ফলাফল পূর্বাভাস করে। Logistic regression এখনো একটি শক্তিশালী ও ব্যাখ্যাযোগ্য বেসলাইন, আর tree
ensemble প্রায়ই চলকের পারস্পরিক ক্রিয়া ভালো ধরতে পারে \citep{kuhn2013}। শ্রেণিগুলো অসম হলে (যেমন এই প্রকল্পে বড় অর্ডার),
accuracy বিভ্রান্তিকর: যে মডেল সবসময় সংখ্যাগরিষ্ঠ শ্রেণি বলে, সেটি নির্ভুল দেখাতে পারে অথচ অকেজো। Precision, recall,
F1-score, ROC বক্ররেখার নিচের ক্ষেত্রফল (ROC-AUC) এবং precision--recall বক্ররেখার নিচের ক্ষেত্রফল বেশি তথ্যবহুল। ভুলের
ব্যবসায়িক ব্যয়ের সঙ্গে মিলিয়ে সিদ্ধান্ত-সীমা (threshold) validation ডেটায় নির্ধারণ করা উচিত \citep{kuhn2013}।

\section{ক্লাস্টারিং ও বিভাজন}
বাজার বিভাজন বিপণন কৌশলের ভিত্তি \citep{kotler2016}। ডেটা-চালিত বিভাজন সাধারণত ক্লাস্টারিংয়ের উপর নির্ভর করে। $K$-means
পর্যবেক্ষণগুলোকে $k$টি দলে ভাগ করে দলের ভেতরের ভেদাঙ্ক ন্যূনতম করে। Agglomerative (Ward) ক্লাস্টারিং ধাপে ধাপে একত্রীকরণের
একটি স্তরক্রম তৈরি করে। Gaussian mixture মডেল প্রতিটি ক্লাস্টারকে একটি বহুচলক normal বণ্টন ধরে নেয় \citep{han2012}। যেহেতু
ক্লাস্টারিংয়ে কোনো ground truth নেই, অভ্যন্তরীণ বৈধতা সূচক ব্যবহৃত হয়, যেমন silhouette coefficient \citep{rousseeuw1987},
Davies--Bouldin সূচক \citep{davies1979} এবং Calinski--Harabasz মানদণ্ড।

\section{অ্যানোমালি ডিটেকশন}
অ্যানোমালি ডিটেকশন এমন পর্যবেক্ষণ শনাক্ত করে যা প্রত্যাশিত আচরণ থেকে বিচ্যুত। কারণ হতে পারে ডেটা-এন্ট্রির ভুল, জালিয়াতি, বা প্রকৃত
কিন্তু ব্যতিক্রমী ঘটনা। প্রচলিত পদ্ধতিগুলো একটি প্রত্যাশিত মান থেকে robust দূরত্ব হিসাব করে। অন্যদিকে Isolation Forest
\citep{liu2008} এলোমেলো পুনরাবৃত্ত বিভাজনের মাধ্যমে বিন্দুগুলোকে বিচ্ছিন্ন করে: অস্বাভাবিক বিন্দু বিচ্ছিন্ন করতে কম বিভাজন লাগে,
ফলে একটি দক্ষ ও বণ্টন-নিরপেক্ষ স্কোর পাওয়া যায়। এটি নিয়ম-ভিত্তিক যাচাইয়ের স্বাভাবিক পরিপূরক।

\section{ব্যাখ্যাযোগ্য মেশিন লার্নিং}
ব্যবস্থাপকরা এমন মডেলে খুব কমই আস্থা রাখেন যা তাঁরা বুঝতে পারেন না। Permutation importance মাপে একটি ফিচার এলোমেলো করলে
মডেলের ত্রুটি কতটা বাড়ে \citep{breiman2001}। সমবায়ী গেম তত্ত্বের উপর ভিত্তি করে SHAP মান প্রতিটি পূর্বাভাসকে যোগাত্মকভাবে
ইনপুট ফিচারগুলোর মধ্যে বণ্টন করে \citep{lundberg2017}। চাহিদা পূর্বাভাস ও বড় অর্ডারের পেছনে কোন উপাদান কাজ করে তা দেখাতে
এখানে দুটিই ব্যবহৃত হয়েছে।

\section{ড্যাশবোর্ড ও দৃশ্যমান সিদ্ধান্ত-সহায়তা}
\citet{few2013} ড্যাশবোর্ডকে সংজ্ঞায়িত করেছেন এক বা একাধিক লক্ষ্য অর্জনে প্রয়োজনীয় সবচেয়ে গুরুত্বপূর্ণ তথ্যের একটি দৃশ্যমান
উপস্থাপনা হিসেবে, যা এক পর্দায় সাজানো যাতে এক নজরে পর্যবেক্ষণ করা যায়। তিনি সাজসজ্জামূলক চার্ট, 3-D প্রভাব, গেজ এবং অর্থহীন
রঙের বিরুদ্ধে সতর্ক করেন; বরং পরিষ্কার তুলনা, সংযত উপাদান ও সামঞ্জস্যপূর্ণ এনকোডিংয়ের পরামর্শ দেন। \cref{ch:dashboard}-এর
ড্যাশবোর্ড নকশায় এই নীতিগুলো অনুসরণ করা হয়েছে। আধুনিক ওপেন-সোর্স টুল দিয়ে বিশ্লেষণের ভাষাতেই সরাসরি ইন্টারঅ্যাকটিভ
ড্যাশবোর্ড তৈরি ও বিনামূল্যে প্রকাশ করা যায়; উদাহরণ: Python লাইব্রেরি pandas \citep{mckinney2022}, scikit-learn
\citep{pedregosa2011}, Plotly ও Streamlit।

\section{গবেষণা-ঘাটতি}
সাহিত্যে প্রতিটি উপাদানের জন্য পরিপক্ব কৌশল রয়েছে: ডেটা পরিষ্কারকরণ, ওয়্যারহাউজিং, পূর্বাভাস, ক্লাসিফিকেশন, ক্লাস্টারিং,
অ্যানোমালি ডিটেকশন ও ড্যাশবোর্ড। কিন্তু গবেষণা ও শিক্ষণ-উপকরণে এগুলো সাধারণত বিচ্ছিন্নভাবে আলোচিত হয়। অনেক প্রায়োগিক পূর্বাভাস
গবেষণায় naive বেঞ্চমার্ক বাদ পড়ে বা এলোমেলো বিভাজনে যাচাই করা হয়, ফলে তাদের নির্ভুলতার দাবি বিশ্বাস করা কঠিন
\citep{makridakis2018, kaufman2012}। খুব কম কাজেই একটি \emph{শুরু থেকে শেষ পর্যন্ত} পুনরুৎপাদনযোগ্য কাঠামো দেখানো হয়েছে,
যেখানে (i) প্রিপ্রসেসিং, মডেলিং ও প্রতিবেদন স্বয়ংক্রিয়ভাবে সংযুক্ত, (ii) মূল্যায়ন কঠোরভাবে সময়-ক্রমিক ও বেঞ্চমার্কসহ, এবং
(iii) ফলাফল কাঁচা মডেল-স্কোর হিসেবে নয়, বরং ব্যাখ্যাসহ পূর্বাভাস ও সুপারিশ হিসেবে ব্যবস্থাপকদের কাছে পৌঁছায়। এই প্রকল্প
বাংলাদেশের একটি ভোগ্যপণ্য বিক্রয় প্রেক্ষাপটে সেই ঘাটতি পূরণ করে, যে বাজার প্রকাশিত অ্যানালিটিক্স সাহিত্যে তুলনামূলক কম প্রতিনিধিত্ব পায়।

\begin{table}[htbp]\centering\small
\caption{মূল সাহিত্যের সারসংক্ষেপ ও এই প্রকল্পে তার ব্যবহার}\label{tab:litsummary}
\begin{tabularx}{\linewidth}{@{}p{4.2cm}X@{}}
\toprule
\textbf{উৎস} & \textbf{এই প্রকল্পে ব্যবহৃত অবদান} \\ \midrule
\citet{shearer2000} & পদ্ধতির ভিত্তি CRISP-DM প্রক্রিয়া মডেল \\
\citet{chen2012}; \citet{sharda2018} & BI\&A-এর বর্ণনামূলক--পূর্বাভাসমূলক--নির্দেশমূলক কাঠামো \\
\citet{garcia2015}; \citet{han2012} & পরিষ্কারকরণ, imputation, outlier ও রূপান্তর কৌশল \\
\citet{kaufman2012} & Leakage-মুক্ত ফিচার ও যাচাই নকশা \\
\citet{kimball2013} & স্টার-স্কিমা ডেটাবেস \\
\citet{hyndman2021}; \citet{winters1960} & Holt--Winters, SARIMA, পূর্বাভাস reconciliation \\
\citet{makridakis2018}; \citet{bergmeir2012} & বেঞ্চমার্ক ও সময়-ক্রমিক যাচাই \\
\citet{breiman2001}; \citet{friedman2001}; \citet{chen2016} & Random forest, gradient boosting, XGBoost \\
\citet{rousseeuw1987}; \citet{davies1979} & ক্লাস্টারের বৈধতা সূচক \\
\citet{liu2008} & Isolation Forest অ্যানোমালি ডিটেকশন \\
\citet{lundberg2017} & SHAP ব্যাখ্যা \\
\citet{few2013} & ড্যাশবোর্ড নকশার নীতি \\ \bottomrule
\end{tabularx}
\end{table}
